% Level 1, Homework 11 -- Get Ready for the Showcase   (input by homework.tex)
\homeworkweek{11}{Get Ready for the Showcase}{15--20 minutes. This week's homework is project work: finish and practise.}
\namedateline

\begin{core}[Finish and practise (about 10 minutes)]
\begin{enumerate}[leftmargin=1.4em,itemsep=2pt]
  \item If your project is missing an ingredient, add it now (or write on paper how you will add it next week). Which ingredients still need work?
  \answerlines{1}
  \item Practise your \textbf{2-minute demo} out loud for someone at home. They are the visitor: let them answer your program's question (or pretend, using your sketch). Use the script in your workbook (chapter 12).
  \item Ask your visitor for \textbf{two stars and a wish}:\par
  \faStar\ \rule{11cm}{0.4pt}\par\smallskip
  \faStar\ \rule{11cm}{0.4pt}\par\smallskip
  \faMagic\ \rule{11cm}{0.4pt}
\end{enumerate}
\end{core}

\begin{challenge}[Be ready for strange answers]
Visitors sometimes type surprising things. Make your program cope:
\begin{itemize}[leftmargin=1.2em,itemsep=1pt]
  \item a number that is too big? Use an \texttt{if} to make it smaller and print a friendly message;
  \item \texttt{Night} instead of \texttt{night}? Add \texttt{.lower()} after \texttt{input(...)}.
\end{itemize}
\end{challenge}

\begin{creative}[Make your stand card]
Fold a piece of paper in half to make a card that stands up. Write your project's title, your \textbf{nickname} as the artist, and one sentence about your artwork. Decorate it! (No real full names, please: nicknames only.)
\end{creative}

\begin{selfcheck}
  \item My program runs from start to end without red errors.
  \item All six ingredients are ticked in the comment list at the top of my file.
  \item I practised my 2-minute demo out loud.
  \item I can point to my function, my loop and my \texttt{if}.
\end{selfcheck}
